% !TeX program = lualatex
% !TeX encoding = utf8
% !TeX spellcheck = uk_UA
% !TeX root =../PyscfBook.tex


%% ========================================================
\section{Метод зв’язаних кластерів}
%% ========================================================


%% --------------------------------------------------------
\subsection{Ідея методу}
%% --------------------------------------------------------

Метод зв'язаних кластерів (\texttt{Coupled Cluster, CC}) є наступним кроком після методу
конфігураційної взаємодії (\texttt{CI}), але з принципово новим підходом
до побудови хвильової функції.

У методі \texttt{CI} хвильова функція системи записується як лінійна комбінація
детермінантів, збуджених відносно основного (\texttt{HF}) детермінанта \eqref{eq:CI}.
Така форма зручна для чисельної реалізації, проте має два
фундаментальні недоліки:
\begin{enumerate}
    \item Щоб отримати точний результат, потрібно включити всі можливі
    збудження (\texttt{Full CI}), що обчислювально неможливо для більшості систем.
    \item Хвильова функція \texttt{CI} не є \textbf{size-extensive}, тобто енергія двох незалежних систем
    не дорівнює сумі їхніх енергій. Це призводить до помилок при описі
    великих молекул.
\end{enumerate}

Метод зв'язаних кластерів (\texttt{CC}) розв'язує ці проблеми, вводячи іншу,
експоненціальну форму хвильової функції:
\begin{equation}\label{eq:CC}
    |\Psi_{CC}\rangle = e^{\hat{T}} |\Phi_0\rangle
\end{equation}
де $|\Phi_0\rangle$ --- HF-детермінант, $\hat{T}$ --- \textbf{кластерний оператор} $\hat{T}$:
\begin{equation}
    \hat{T} = \hat{T}_1 + \hat{T}_2 + \hat{T}_3 + \cdots,
\end{equation}
а кожен член \(\hat{T}_n\) відповідає всім можливим \(n\)-кратним збудженням
відносно детермінанта Хартрі–Фока $|\Phi_0\rangle$:
\begin{align}
    \hat{T}_1 &= \sum_{i,a} t_i^a\, a_a^\dagger a_i, \\
    \hat{T}_2 &= \frac{1}{4} \sum_{i,j,a,b} t_{ij}^{ab}\, a_a^\dagger a_b^\dagger a_j a_i, \\
    &\vdots
\end{align}

Коефіцієнти (кластерні амплітуди) $t$ є невідомими параметрами, які слід визначити.

\begin{commentbox}
На перший погляд, експоненціальна форма виглядає складнішою, ніж
звичайна лінійна комбінація. Проте її фундаментальна перевага полягає
у властивості експоненти:
\begin{equation}
    e^{\hat{T}} = 1 + \hat{T} + \frac{1}{2!}\hat{T}^2 +
    \frac{1}{3!}\hat{T}^3 + \cdots,
\end{equation}
що означає, що навіть якщо явно враховуються лише \(\hat{T}_1\) і \(\hat{T}_2\)
(метод \texttt{CCSD}), хвильова функція автоматично містить внески від
трикратних, чотирикратних та вищих збуджень через добутки
\(\hat{T}_1^2\), \(\hat{T}_1\hat{T}_2\) тощо. Тобто експонента  $e^{\hat{T}}$ автоматично породжує всі можливі добутки збуджень. Тому метод \texttt{CC} враховує складені конфігурації навіть при обмеженій кількості параметрів ---
на відміну від лінійного \texttt{CI}.
\end{commentbox}

\paragraph{Отримання рівнянь для кластерних амплітуд.}
Підставимо анзац \eqref{eq:CC} у рівняння Шредінґера і домножимо зліва на $e^{-\hat{T}}$:
\[
\bar{H}\,|\Phi_0\rangle = E_{\text{CC}}\,|\Phi_0\rangle, \qquad
\bar{H} = e^{-\hat{T}}\hat{H}e^{\hat{T}}.
\]
Тоді енергія та амплітуди визначаються як:
\[
E_{\text{CC}} = \langle \Phi_0 | \bar{H} | \Phi_0 \rangle, \qquad
\langle \Phi_\mu | \bar{H} | \Phi_0 \rangle = 0, \quad
(\mu = \text{одно-, дво- та ін. збудження}).
\]
Ці рівняння є нелінійними відносно кластерних амплітуд \(t\),
оскільки оператор $e^{\hat{T}}$ містить добутки $t_i^a,\,t_{ij}^{ab},\,\ldots$.
Тому вони розв’язуються \textbf{ітераційно} --- починаючи з початкового наближення (наприклад, MP2),
доки енергія $E_{\text{CC}}$ не мінімізується.

\begin{commentbox}
Умова
\[
\langle \Phi_\mu | \bar{H} | \Phi_0 \rangle = 0
\]
означає, що після врахування всіх збуджень через оператор \(\hat{T}\),
ефективний гамільтоніан
\(\bar{H} = e^{-\hat{T}}\hat{H}e^{\hat{T}}\)
більше не створює нових збуджень із основного стану \(|\Phi_0\rangle\).
Тобто амплітуди \(t\) підбираються так, щоб усі збудження були вже враховані
в експоненті, і хвильова функція \(e^{\hat{T}}|\Phi_0\rangle\) стала власним
станом гамільтоніана.
\end{commentbox}

%%% --------------------------------------------------------
%\subsection{Доведення розмірної екстенсивності.}
%%% --------------------------------------------------------
%
%Для двох незалежних підсистем $A$ і $B$ маємо
%$\hat{H} = \hat{H}_A + \hat{H}_B$,
%$\Phi_0 = \Phi_0^A \otimes \Phi_0^B$,
%$\hat{T} = \hat{T}_A + \hat{T}_B$.
%Оскільки оператори діють у різних просторах:
%\[
%e^{\hat{T}} = e^{\hat{T}_A}e^{\hat{T}_B}, \quad
%\bar{H} = e^{-\hat{T}}\hat{H}e^{\hat{T}} =
%e^{-\hat{T}_A}\hat{H}_A e^{\hat{T}_A} +
%e^{-\hat{T}_B}\hat{H}_B e^{\hat{T}_B}.
%\]
%Звідси випливає адитивність енергії:
%\[
%E_{\text{CC}} = E_{\text{CC}}(A) + E_{\text{CC}}(B),
%\]
%тобто метод \texttt{CC} є \emph{size-extensive} — правильно масштабується з розміром системи.

%% --------------------------------------------------------
\subsection{Наближення в методі CC: \texttt{CCSD} і \texttt{CCSD(T)}}
%% --------------------------------------------------------

У практичних розрахунках повна експоненційна форма хвильової функції методу
зв'язаних кластерів:
\[
|\Psi_{\text{CC}}\rangle = e^{\hat{T}} |\Phi_0\rangle,
\]
де
\[
\hat{T} = \hat{T}_1 + \hat{T}_2 + \hat{T}_3 + \dots + \hat{T}_N,
\]
залишається, звісно, нескінченною, адже кількість можливих багатоелектронних
збуджень зростає комбінаційно з числом електронів.
Тому необхідно робити \textbf{наближення}, обмежуючи порядок збуджень.

\paragraph{Наближення \texttt{CCSD}}
Найпоширеніше робоче наближення — це
\[
\hat{T} \approx \hat{T}_1 + \hat{T}_2,
\]
тобто враховуються лише \textbf{однократні} та \textbf{дворазові збудження}.
Таке наближення позначається як \texttt{CCSD}
(від англ. \emph{Coupled Cluster with Single and Double excitations}).

Оператор хвильової функції має вигляд:
\[
|\Psi_{\text{CCSD}}\rangle = e^{\hat{T}_1 + \hat{T}_2} |\Phi_0\rangle.
\]

Хоча формально ми не включаємо потрійні й вищі збудження,
експоненційна форма гарантує, що частина їхнього ефекту враховується автоматично
через добутки нижчих операторів, наприклад:
\[
\frac{1}{2!}\hat{T}_1^2 \sim \text{подвійні збудження}, \quad
\hat{T}_1 \hat{T}_2 \sim \text{потрійні збудження}.
\]
Тобто навіть \texttt{CCSD} описує взаємодію між різними рівнями кореляції значно точніше,
ніж метод \texttt{CISD} (конфігураційна взаємодія з одиничними та подвійними збудженнями).

\paragraph{Корекція \texttt{(T)}: «золотий стандарт»}
Для більшої точності часто додають внесок потрійних збуджень \(\hat{T}_3\)
у вигляді \emph{неітераційної поправки}, яку обчислюють на основі вже отриманих
амплітуд \(\hat{T}_1, \hat{T}_2\).
Таке наближення позначається як \texttt{CCSD(T)}.

Символ \texttt{(T)} означає, що потрійні збудження враховані
\emph{perturbatively}, тобто за теорією збурень.
Така схема зберігає високу точність, наближену до повного
\texttt{CCSDT}, але з обчислювальною складністю всього
\(\mathcal{O}(N^7)\) замість \(\mathcal{O}(N^8)\)–\(\mathcal{O}(N^{10})\)
для повного варіанта.

\paragraph{Практичне значення}
Метод \texttt{CCSD(T)} виявився настільки надійним у відтворенні
експериментальних енергій, геометрій і частот коливань, що його часто називають
\textbf{«золотим стандартом» квантової хімії}.
Для невеликих і середніх молекул похибка енергії зв’язку зазвичай не перевищує
\(1{-}2~\text{kcal/mol}\), тобто досягає хімічної точності.

\begin{center}
\begin{tblr}{
colspec={Q[l,m] Q[l,m] Q[l,m]},
row{1} = {c,m, font=\bfseries}
}
\toprule
Метод & Враховані збудження & Тип наближення \\ \midrule
\texttt{CCSD} & \(\hat{T}_1, \hat{T}_2\) & Ітераційне \\
\texttt{CCSD(T)} & \(\hat{T}_1, \hat{T}_2\) + корекція \(\hat{T}_3\) & Ітераційно + збурювально \\
\texttt{CCSDT} & \(\hat{T}_1, \hat{T}_2, \hat{T}_3\) & Повне, ітераційне \\ \bottomrule
\end{tblr}
\end{center}

Таким чином, \texttt{CCSD(T)} поєднує фізичну повноту опису з практичною
ефективністю й залишається основним методом високоточної
квантової хімії для систем до кількох десятків атомів.

%% --------------------------------------------------------
\subsection{Приклади розрахунків у PySCF}
%% --------------------------------------------------------

\subsubsection{Базове використання методів CC}

Розглянемо послідовно, як увімкнути різні варіанти методу зв'язаних кластерів
у бібліотеці PySCF:

\begin{minted}{python}
from pyscf import gto, scf, cc

mol = gto.M(
    atom = 'H 0 0 0; F 0 0 1.1',
    basis = 'cc-pvqz',
    verbose = 0
)

# Розрахунок HF
mf = scf.RHF(mol).run()

# CCSD
cc = cc.CCSD(mf)
cc.kernel()

# CCSD з наступною (T)-корекцією
et = mycc.ccsd_t()

print(f"Енергія HF: {mf.e_tot:.4f} Ha")
print(f"Енергія CCSD: {cc.e_tot:.4f} Ha")
print(f"Поправка (T): {et:.4f} Ha")
print(f"Енергія CCSD(T): {cc.e_tot + et:.4f} Ha")
print(f"Експеримент: {-100.4489} Ha")
\end{minted}

